% !PRINTED 77
\weberpage{77}

%<<CONT
not contained in $\mathfrak{B}'$, and hence, if $\beta_0', \beta_1', \dots \beta_\mu'$ are the
indices of $b'$, then for every $x$ one of the congruences

\weq{16}{
\begin{aligned}
x\beta_0' + \beta_0 &\equiv 0 \pmod{2} \\
x\beta_1' + \beta_1 &\equiv 0 \pmod{q_1} \\
&\cdots\cdots\cdots\cdots \\
x\beta_\mu' + \beta_\mu &\equiv 0 \pmod{q_\mu}
\end{aligned}
}

must be satisfied.

If now two numbers $\beta', \beta$ are not both divisible by a prime $q$, then one can always dispose of $x$ in such a way that $x\beta' + \beta$ does not become divisible by $q$; one need only assume $x$ not divisible by $q$ when $\beta$ is divisible by $q$, and $x$ divisible by $q$ when $\beta$ is not divisible by $q$. One can therefore also determine $x$ so that it satisfies several such requirements simultaneously. It therefore follows from (16) that either $\beta_0', \beta_0$ must be divisible by $2$, or $\beta_1', \beta_1$ by $q_1 \dots$ or $\beta_\mu', \beta_\mu$ by $q_\mu$; i.\,e. one of the congruences (15) must also be fulfilled when $b$ is any number from $\mathfrak{B}$, no matter whether it occurs in $\mathfrak{B}'$ or not. Finally it follows again from (16), if we understand by $b, b'$ two arbitrary numbers from $\mathfrak{B}$, that one and the same of the congruences (15) must hold for all numbers from $\mathfrak{B}$. We can now summarise this as follows:

\gesperrt{2. If the period $\eta$ does not belong to $\mathfrak{A}$, then among the primes (including 2) that divide $m$ more than once there must be one, $q$, such that for every number $b$ from $\mathfrak{B}$ the index $\beta$ corresponding to $q$ is divisible by $q$.}

We now consider the group $Q$ which, according to the notation (5) of \S.\,18, is defined by the congruence

\weq{17}{Q \equiv 1 \pmod{\dfrac{m}{q}}}

and denote by $P$ the group reciprocal to it.

By Theorem 7., \S.\,19, every group which, like $\mathfrak{B}$, has the property that the index $\beta$ corresponding to $q$ is divisible by $q$, is a divisor of the group $P$, and consequently, by \S.\,12, 8., $Q$ is a divisor of $\mathfrak{A}$; hence, by \S.\,19, 8., $\mathfrak{A}$ is a non-primary divisor of $\mathfrak{N}$, and with this we have arrived at the theorem:

\gesperrt{3. If $\mathfrak{A}$ is a primary divisor of $\mathfrak{N}$, then the period $\eta$ belongs to the group $\mathfrak{A}$.}
